\documentclass[11pt]{amsart}

\usepackage[T1]{fontenc}
\usepackage{lmodern}
\usepackage{microtype}
\usepackage{mathtools,amssymb,amsthm}
\usepackage{booktabs}
\usepackage[hidelinks]{hyperref}

\hypersetup{
  pdftitle={A Nine-Element Counterexample to Sinclair's Front-Loading Conjecture}
}

\newtheorem{theorem}{Theorem}
\theoremstyle{remark}
\newtheorem{remark}[theorem]{Remark}

\newcommand{\rk}{\operatorname{rk}}
\newcommand{\cl}{\operatorname{cl}}

\title[A counterexample to front-loading]
{A Nine-Element Counterexample to Sinclair's Front-Loading Conjecture}

\date{}

\subjclass[2020]{05B35}
\keywords{matroid, graded M\"obius algebra, differential Hilbert function,
counterexample}

\begin{document}

\begin{abstract}
Sinclair's Conjecture~C asserts that the defect of every simple matroid is
front-loaded. We give a simple rank-five matroid $M$ on nine elements,
representable over $\mathbb Q$, whose apolar Hilbert series and
radial-derivative Hilbert series, both recalled in Section~\ref{sec:ce}, are
\[
H_M^0(q)=1+9q+27q^2+27q^3+9q^4+q^5
\]
and
\[
H_{\beta,M}(q)=1+9q+31q^2+32q^3+9q^4+q^5.
\]
Here $h_k^\beta(M)$ and $h_k^0(M)$ denote the coefficients of $q^k$ in
$H_{\beta,M}$ and in $H_M^0$, so the defect $\delta_k(M)=h_k^\beta(M)-h_k^0(M)$
has $\delta_2(M)=4<5=\delta_3(M)$, which refutes the conjecture. The calculation
uses exact rational arithmetic and is reproduced by the short
standard-library program printed in the appendix.
\end{abstract}

\maketitle

\section{The counterexample}\label{sec:ce}

We use Sinclair's notation~\cite{Sinclair}. Let $M$ be a simple matroid of
rank $r$ on a ground set $E$, and let $g_E$ be the top class of its graded
M\"obius algebra, rescaled as in the proof of
Theorem~\ref{thm:counterexample}: there $g_X=\beta_Xe_X$ for a flat $X$, where
$\beta_X$ is the number of ordered bases of $X$ and $e_X$ is the corresponding
basis vector of the graded M\"obius algebra. Write
\[
H_{\beta,M}(q)=\sum_{k=0}^r h_k^\beta(M)q^k
\]
for the Hilbert series of the cyclic module generated by $g_E$ under the
radial derivative $D_\beta$ and the coordinate derivatives $L_a^\beta$
($a\in E$), graded by the number of operator applications: in the notation of
the chain \eqref{eq:chain} below, $h_k^\beta(M)=\dim_{\mathbb Q}\mathcal
C_k^\beta$. Omitting $D_\beta$ gives the classical apolar Hilbert function
$h_k^0(M)=\dim_{\mathbb Q}\mathcal C_k^0$, and we write
\[
H_M^0(q)=\sum_{k=0}^r h_k^0(M)q^k.
\]
Sinclair's Conjecture~C states that, with
\[
\delta_k(M)=h_k^\beta(M)-h_k^0(M),
\]
one has
\[
\delta_k(M)\geq \delta_{r-k}(M),\qquad k\leq r/2.
\]
Sinclair records the equivalent form $h_k^\beta(M)\geq h_{r-k}^\beta(M)$ for
$k\leq r/2$~\cite{Sinclair}; the two forms agree for the matroid of
Theorem~\ref{thm:counterexample}, whose $h_k^0$ is symmetric under $k\mapsto
r-k$. The introduction of~\cite{Sinclair} describes the same inequality as
top-heaviness of $H_\beta$ in differential degree.

\begin{theorem}\label{thm:counterexample}
Let $M$ be the column matroid over $\mathbb Q$ of
\[
A=
\begin{pmatrix}
1&0&0&0&0&1&0&0&0\\
0&1&0&0&0&0&1&1&0\\
0&0&1&0&0&1&1&2&0\\
0&0&0&1&0&0&0&0&1\\
0&0&0&0&1&0&1&3&1
\end{pmatrix}.
\]
Then $M$ is a simple rank-five matroid on nine elements and
\begin{align*}
H_M^0(q)
 &=1+9q+27q^2+27q^3+9q^4+q^5,\\
H_{\beta,M}(q)
 &=1+9q+31q^2+32q^3+9q^4+q^5.
\end{align*}
Consequently, $M$ is a counterexample to Conjecture~C of
Sinclair~\cite{Sinclair}.
\end{theorem}

\begin{proof}
The first five columns of $A$ form the identity matrix, so $\rk(M)=5$.
Every column is nonzero and no two columns are proportional; hence $M$ is
simple.

For a flat $X$, let $\beta_X$ be the number of ordered bases of $X$, let
$e_X$ denote the corresponding basis vector of the graded M\"obius algebra,
and put $g_X=\beta_Xe_X$. By Sinclair's operator formula
\cite[Proposition~2.4]{Sinclair},
\begin{equation}\label{eq:operators}
L_a^\beta g_Y
 =\sum_{\substack{X\lessdot Y\\X\vee a=Y}}g_X,
\qquad
D_\beta g_Y
 =\rk(Y)\sum_{X\lessdot Y}
   \frac{|Y\setminus X|}{9-|X|}\,g_X.
\end{equation}
Indeed, the atoms are the elements of $E$, and for a cover
$X\lessdot Y$ one has $X\vee a=Y$ exactly when $a\in Y\setminus X$.
Thus the quantities in Sinclair's formula satisfy
\[
s(X)=9-|X|,
\qquad
m(X,Y)=|Y\setminus X|.
\]
Both are evaluated at the lower flat of the cover: the sets $Y\setminus X$,
as $Y$ runs over the covers of $X$, partition the atoms not lying below $X$,
so that $\sum_{Y\gtrdot X}m(X,Y)=9-|X|=s(X)$.

Starting from $\mathcal C_0^\beta=\mathbb Qg_E$, define
\begin{equation}\label{eq:chain}
\mathcal C_{k+1}^\beta
 =D_\beta\mathcal C_k^\beta
  +\sum_{a\in E}L_a^\beta\mathcal C_k^\beta,
\end{equation}
and define $\mathcal C_k^0$ by omitting $D_\beta$; as in
Section~\ref{sec:ce}, $h_k^0(M)=\dim_{\mathbb Q}\mathcal C_k^0$ and
$h_k^\beta(M)=\dim_{\mathbb Q}\mathcal C_k^\beta$, which are the numbers
tabulated below. All matrices in \eqref{eq:operators} have rational entries,
so their ranks over $\mathbb Q$ agree with their ranks over $\mathbb R$.

For each $S\subseteq E$, exact elimination computes $\rk(A_S)$ and hence
\[
\cl(S)=\{a\in E:\rk(A_{S\cup\{a\}})=\rk(A_S)\}.
\]
Enumerating the $2^9$ subsets gives the Whitney numbers
\[
(W_0,W_1,W_2,W_3,W_4,W_5)=(1,9,32,49,27,1).
\]
Constructing the flat-cover matrices in \eqref{eq:operators} and applying
exact Gauss--Jordan elimination degree by degree gives
\[
\begin{array}{c@{\qquad}rrrrrr}
\toprule
k&0&1&2&3&4&5\\
\midrule
W_{5-k}(M)&1&27&49&32&9&1\\
h_k^0(M)&1&9&27&27&9&1\\
h_k^\beta(M)&1&9&31&32&9&1\\
\delta_k(M)&0&0&4&5&0&0\\
\bottomrule
\end{array}
\]
For $r=5$, the index $k=2$ lies in the conjectured range, but
\[
h_2^\beta(M)=31<32=h_3^\beta(M),
\qquad
\delta_2(M)=4<5=\delta_3(M).
\]
This refutes Conjecture~C.
\end{proof}

\begin{remark}
The sequence $(1,9,31,32,9,1)$ is strictly log-concave, so the example does
not contradict Sinclair's separate log-concavity conjecture. No claim of
minimality is made.
\end{remark}

\paragraph{Status.}
Conjecture~C fails for $M$ in both of the forms printed
in~\cite{Sinclair}, and only in the range $k\leq r/2$ that the conjecture
addresses; the cases proved there --- the endpoints $\delta_0=\delta_r=0$ and
the uniform matroids, where $H_\beta=H^0$ --- are untouched, and $M$ is
neither. As we read~\cite{Sinclair}, Conjecture~C is verified there by exact
rational elimination for all $950$ pairwise nonisomorphic simple matroids on
eight elements \cite[Prop.~5.1]{Sinclair}, and a modular check of a stratified
sample of $400$ simple nine-element matroids is reported
\cite[Sec.~5]{Sinclair} and called there evidence rather than an exact
certificate; that sample is specified by a seed, and we make no claim about
whether $M$ lies in it. We have not reverified either of those computations,
and Theorem~\ref{thm:counterexample} does not depend on them: they bear only
on where a counterexample could already have been found. On the same reading,
no exhaustive nine-element census is asserted in~\cite{Sinclair}, and we are
not aware of any other work that resolves Conjecture~C or exhibits a
counterexample to it; the invariants $H_\beta$, $D_\beta$ and $\delta_k$ are
introduced in~\cite{Sinclair}. Every statement about \cite{Sinclair} in this
paragraph is a reading of the version cited below, not an output of the
computation reported next.

\section*{Exact verification}

The computation is a direct enumeration in exact rational arithmetic. All
$2^9$ subsets of $E$ are ranked by elimination over $\mathbb Q$, which yields
the closure operator, the $119$ flats, the Whitney numbers and the covering
relations; the matrices of \eqref{eq:operators} are then assembled on the
basis $\{g_X\}$ and the chain \eqref{eq:chain} is reduced one degree at a
time. The program in Appendix~\ref{app:code} carries this out with the Python
standard library alone, every arithmetic operation performed on
\texttt{Fraction} objects, and prints the rank and simplicity of $M$, its
number of bases ($67$), its number of flats ($119$), the Whitney numbers,
the join step used in the proof above, both Hilbert functions and $\delta$.
No floating-point computation and no external matroid catalogue is used.

The same assertions were recomputed independently on a separate machine, by a
longer program written for this check from the statements above rather than
from the listing below, again using the Python standard library alone and
exact arithmetic throughout. It reports $72$ checks and all $72$ pass; it
agrees with every number printed here, namely rank five, simplicity, $67$
bases, $119$ flats, $(1,9,32,49,27,1)$, $(1,9,27,27,9,1)$, $(1,9,31,32,9,1)$
and $(0,0,4,5,0,0)$, and it brackets each dimension from both sides, by a
nonvanishing minor of that size and by explicit annihilating functionals. As a
control it was run on the uniform matroid $U_{5,9}$, which is determined up to
isomorphism by its rank and its number of elements, and there it reports that
Conjecture~C holds; so the test is not one that fires on every input. The
first program is printed in full in Appendix~\ref{app:code}, and the longer
program together with the transcript of its run accompanies this note as
supporting material.

Two programs agreeing is not a proof, and two things lie outside the reach of
both, each being a matter of reading~\cite{Sinclair} rather than of
computation: the transcription of the operator formula \eqref{eq:operators},
and the convention that $\delta_k$ is indexed by the number of operator
applications, as in \eqref{eq:chain} and in the table above. Both are printed
in full above rather than only cited --- \eqref{eq:operators} together with the
identifications $s(X)=9-|X|$ and $m(X,Y)=|Y\setminus X|$ that specialize it to
$M$, and the grading in \eqref{eq:chain} --- so that they can be compared with
the source directly. The Hilbert functions of
Theorem~\ref{thm:counterexample} are exact consequences of the definitions of
Section~\ref{sec:ce} and do not depend on that comparison; the identification
of the inequality they violate with Conjecture~C does. The grading is the
load-bearing half: indexing instead by the degree of the graded piece,
$h_j=\dim_{\mathbb Q}\mathcal C_{r-j}$, replaces $\delta$ by $(0,0,5,4,0,0)$,
which satisfies the inequality. There is no census or minimality claim here,
and hence no exhaustive search to reproduce.

\raggedbottom

\appendix

\section{The verification program}\label{app:code}

The following program is self-contained and recomputes every number in the
table of Theorem~\ref{thm:counterexample} from the matrix $A$ alone.
Comments give the expected output, one \texttt{print} statement per line.

{\small
\begin{verbatim}
from fractions import Fraction as F
from itertools import combinations
from math import factorial

A = [[1, 0, 0, 0, 0, 1, 0, 0, 0],
     [0, 1, 0, 0, 0, 0, 1, 1, 0],
     [0, 0, 1, 0, 0, 1, 1, 2, 0],
     [0, 0, 0, 1, 0, 0, 0, 0, 1],
     [0, 0, 0, 0, 1, 0, 1, 3, 1]]
n, r = 9, 5
E = frozenset(range(n))

def rref(rows):                   # exact Gauss-Jordan; returns
    M = [list(v) for v in rows]   # the nonzero reduced rows
    p = 0
    for c in range(len(M[0])):
        piv = next((i for i in range(p, len(M))
                    if M[i][c] != 0), None)
        if piv is None:
            continue
        M[p], M[piv] = M[piv], M[p]
        M[p] = [x / M[p][c] for x in M[p]]
        for i in range(len(M)):
            if i != p and M[i][c] != 0:
                f = M[i][c]
                M[i] = [x - f * y for x, y in zip(M[i], M[p])]
        p += 1
    return M[:p]

def rank_cols(S):
    S = sorted(S)
    return len(rref([[F(A[i][c]) for c in S]
                     for i in range(r)])) if S else 0

rk = {frozenset(S): rank_cols(S)
      for k in range(n + 1) for S in combinations(range(n), k)}
print(rk[E],
      [a for a in range(n) if rk[frozenset([a])] == 0],
      [p for p in combinations(range(n), 2)
       if rk[frozenset(p)] == 1])
#      5 [] []                      rank five and simple
print(sum(1 for S in combinations(range(n), r)
          if rk[frozenset(S)] == r))
#      67                           bases

def cl(S):
    S = frozenset(S)
    return frozenset(a for a in range(n) if rk[S | {a}] == rk[S])

flats = sorted({cl(S) for S in rk},
               key=lambda X: (rk[X], sorted(X)))
lev = {k: [X for X in flats if rk[X] == k] for k in range(r + 1)}
pos = {X: i for k in lev for i, X in enumerate(lev[k])}
print(len(flats), [len(lev[k]) for k in range(r + 1)])
#      119 [1, 9, 32, 49, 27, 1]    flats, then the Whitney numbers

und = {Y: [X for X in lev[rk[Y] - 1] if X < Y]   # covers X < Y
       for k in range(1, r + 1) for Y in lev[k]}
print(all((cl(X | {a}) == Y) == (a in Y - X)
          for Y in und for X in und[Y] for a in range(n)))
#      True                         the join step used in the proof

def beta(X):                    # ordered bases of the flat X
    d = rk[X]                   # beta(E) = 8040 = 67 * 5!
    return factorial(d) * sum(1 for S in combinations(sorted(X), d)
                              if rk[frozenset(S)] == d)

def L(v, k, a):                 # v is indexed by lev[k]
    out = [F(0)] * len(lev[k - 1])
    for Y, c in zip(lev[k], v):
        for X in (und[Y] if c else []):
            if cl(X | {a}) == Y:
                out[pos[X]] += c
    return out

def D(v, k):
    out = [F(0)] * len(lev[k - 1])
    for Y, c in zip(lev[k], v):
        for X in (und[Y] if c else []):
            out[pos[X]] += c * rk[Y] * F(len(Y - X), n - len(X))
    return out

def hilb(use_D):
    v, k = [[F(beta(E))]], r    # C_0 = Q g_E
    h = [len(rref(v))]
    while k:
        nxt = ([D(u, k) for u in v] if use_D else [])
        nxt += [L(u, k, a) for u in v for a in range(n)]
        v = rref(nxt)
        h.append(len(v))
        k -= 1
    return h

h0, hb = hilb(False), hilb(True)
print(h0, hb, [b - a for a, b in zip(h0, hb)])
#      [1, 9, 27, 27, 9, 1] [1, 9, 31, 32, 9, 1] [0, 0, 4, 5, 0, 0]
print(hb[2] < hb[3], hb[2] - h0[2] < hb[3] - h0[3], 2 <= r / 2)
#      True True True               Conjecture C fails at k = 2
\end{verbatim}
}

\begin{thebibliography}{1}

\bibitem{Sinclair}
T.~Sinclair,
\emph{The radial derivative on the graded M\"obius algebra},
arXiv:2608.18519v1 [math.CO], 2026.

\end{thebibliography}

\end{document}
